\documentclass[12pt]{article}
\usepackage[brazil]{babel}
\usepackage[utf8]{inputenc}
\usepackage[T1]{fontenc}
\usepackage{amsmath, amssymb}
\usepackage{geometry}
\usepackage{graphicx}
\usepackage{xcolor}

\definecolor{univBlue}{RGB}{0, 57, 145}
\definecolor{univLightBlue}{RGB}{71, 115, 169}
\definecolor{univSoftBlue}{RGB}{147, 172, 203}
\definecolor{maincolor}{RGB}{0, 57, 145}

\usepackage{fancyhdr}
\usepackage[most]{tcolorbox}

\newtcolorbox{block}[1]{%
  colback=univSoftBlue!10,
  colframe=maincolor,
  coltitle=white,
  title=#1,
  fonttitle=\bfseries,
  boxrule=1.5pt,
  arc=3mm,
  left=6pt,
  right=6pt,
  top=6pt,
  bottom=6pt
}

\geometry{margin=2.5cm}
\pagestyle{fancy}
\fancyhf{}
\usepackage{tikz}
\usetikzlibrary{positioning,arrows.meta,shapes.multipart}

\renewcommand{\headrulewidth}{2pt}
\renewcommand{\footrulewidth}{2pt}

\renewcommand{\headrule}{\hbox to\headwidth{%
  \color{maincolor}\leaders\hrule height \headrulewidth\hfill}}

\renewcommand{\footrule}{\hbox to\headwidth{%
  \color{maincolor}\leaders\hrule height \footrulewidth\hfill}}

\setlength{\headheight}{52pt}

\fancyhead[L]{%
  \includegraphics[height=1.2cm]{universidade_logo_novo.png}
}

\fancyhead[R]{%
  \raggedleft
  \textcolor{maincolor}{%
    \small
    \textbf{Lista de Exercícios — Interfaces e Collections}\\
    Me.~Sergio Souza Novak\\
    \textit{Faculdade de Engenharia de Software}
  }
}

\fancyfoot[C]{%
  \textcolor{maincolor}{\thepage}
}

% Estilo reutilizável para diagramas UML
\tikzset{
  uml/.style={
    draw=maincolor,
    rectangle split,
    rectangle split parts=3,
    fill=univSoftBlue!20,
    minimum width=5cm,
    font=\small\ttfamily,
    align=left
  },
  heranca/.style={-{Triangle[open]}, thick, maincolor},
  implementacao/.style={-{Triangle[open]}, dashed, thick, maincolor},
  composicao/.style={-{Diamond[fill=maincolor]}, thick, maincolor},
  dependencia/.style={-stealth, dashed, thick, maincolor}
}

\begin{document}
% ============================================================
Lista de implementação abordando \textbf{Interfaces} e o \textbf{Java Collections Framework}. Implemente na linguagem Java todos os exercícios propostos. O objetivo é aplicar os conceitos vistos em aula: definição e implementação de interfaces, polimorfismo, e o uso de \texttt{List} (\texttt{ArrayList}, \texttt{LinkedList}), \texttt{Set} (\texttt{HashSet}, \texttt{TreeSet}), \texttt{Map} (\texttt{HashMap}) e ordenação (\texttt{Collections.sort}).

\section*{Formato de Entrega}

\begin{block}{Instruções para Submissão}
A entrega desta lista de exercícios deve ser feita através de um \textbf{único arquivo .zip}.
Dentro deste arquivo, deve haver uma pasta estruturada para cada exercício (ex: \texttt{Exercicio1}, \texttt{Exercicio2}, ..., \texttt{Exercicio6}).

\textbf{Atenção:} Certifique-se de que cada pasta contenha \textbf{todos} os arquivos necessários (\texttt{.java}) para a correta compilação e execução do respectivo exercício.
\end{block}

% ============================================================
\section*{1) Implementação de Interface e Polimorfismo}
Modele um sistema simples de autenticação que demonstra o uso de interfaces e polimorfismo.

\begin{center}
\begin{tikzpicture}[node distance=2cm and 1cm, , scale=0.8, transform shape]

\node[uml, rectangle split parts=2, fill=univSoftBlue!40] (autenticavel) {
    \textbf{\guillemotleft interface\guillemotright} \\
    \textbf{Autenticavel}
    \nodepart{second}
    + autenticar(senha: String) : boolean \\
    + fazerLogout() : void
};

\node[uml, below left=1.5cm and -1cm of autenticavel] (gerente) {
    \textbf{Gerente}
    \nodepart{second}
    - senhaDoGerente : String
    \nodepart{third}
    + autenticar(senha: String) : boolean \\
    + fazerLogout() : void
};

\node[uml, below right=1.5cm and -1cm of autenticavel] (cliente) {
    \textbf{Cliente}
    \nodepart{second}
    - tokenApp : String
    \nodepart{third}
    + autenticar(senha: String) : boolean \\
    + fazerLogout() : void
};

\draw[implementacao] (gerente.north) -- ++(0,0.5) -| (autenticavel.south);
\draw[implementacao] (cliente.north) -- ++(0,0.5) -| (autenticavel.south);

\end{tikzpicture}
\end{center}

\subsection*{Requisitos}
\begin{itemize}
    \item Crie a interface \texttt{Autenticavel} com os métodos indicados no diagrama.
    \item Crie as classes \texttt{Gerente} e \texttt{Cliente} implementando a interface. Cada classe deve ter lógicas diferentes de autenticação (ex: o Gerente compara uma string fixa, o Cliente valida um token simulado).
    \item Crie uma classe \texttt{SistemaLogin} com um método \texttt{fazerLogin(Autenticavel usuario, String senha)} que recebe \textbf{qualquer} objeto que assine o contrato da interface.
    \item No \texttt{Main}, instancie gerentes e clientes, e passe-os para o método \texttt{fazerLogin} do sistema.
\end{itemize}

% ============================================================
\section*{2) ArrayList e Manipulação de Elementos}
Desenvolva um pequeno sistema orientado a objetos para gerenciar um carrinho de compras de um supermercado, utilizando a estrutura de dados \texttt{ArrayList}.

\begin{center}
\begin{tikzpicture}[node distance=1cm and 1cm, scale=0.8, transform shape]
\scriptsize  
\node[uml] (item) {
    \textbf{ItemCompra}
    \nodepart{second}
    - nome : String \\
    - quantidade : int \\
    - precoUnitario : double
    \nodepart{third}
    + ItemCompra(nome: String, qtd: int, preco: double) \\
    + getNome() : String \\
    + calcularSubtotal() : double \\
    + toString() : String
};

\node[uml, right=of item] (carrinho) {
    \textbf{CarrinhoDeCompras}
    \nodepart{second}
    - itens : List<ItemCompra>
    \nodepart{third}
    + adicionarItem(item: ItemCompra) : void \\
    + removerItem(nome: String) : void \\
    + listarItens() : void \\
    + calcularTotal() : double
};

\draw[composicao] (carrinho.west) -- (item.east);

\end{tikzpicture}
\end{center}

\subsection*{Requisitos}
\begin{itemize}
    \item Crie a classe \texttt{ItemCompra} devidamente encapsulada. O método \texttt{calcularSubtotal()} deve retornar a multiplicação da quantidade pelo preço unitário.
    \item Na classe \texttt{CarrinhoDeCompras}, o atributo \texttt{itens} deve ser declarado utilizando a interface \texttt{List} e instanciado como um \texttt{ArrayList}.
    \item Implemente \texttt{adicionarItem} utilizando o método \texttt{add()} da coleção.
    \item Implemente \texttt{removerItem} para buscar o item pelo nome e removê-lo. (Dica: pesquise sobre o método \texttt{removeIf()} da interface \texttt{Collection} ou utilize um \texttt{Iterator}).
    \item O método \texttt{listarItens} deve usar o laço \texttt{for-each} para percorrer a lista e exibir o \texttt{toString()} e o subtotal de cada item. O método \texttt{calcularTotal()} deve somar os subtotais de todos os itens do carrinho.
    \item Crie uma classe principal (\texttt{Main}) para instanciar o carrinho, adicionar alguns itens, remover pelo menos um deles e exibir o resumo final e o valor total.
\end{itemize}

% ============================================================
\section*{3) LinkedList: Fila de Atendimento}
Simule uma fila de atendimento de um hospital, onde a ordem de chegada dita a ordem de atendimento, mas emergências podem furar a fila.

\begin{center}
\begin{tikzpicture}[node distance=2cm and 1cm]

\node[uml] (paciente) {
    \textbf{Paciente}
    \nodepart{second}
    - nome : String \\
    - prioridade : int
    \nodepart{third}
    + getNome() : String \\
    + getPrioridade() : int
};

\node[uml, right=of paciente] (fila) {
    \textbf{GerenciadorFila}
    \nodepart{second}
    - fila : LinkedList<Paciente>
    \nodepart{third}
    + adicionarPaciente(p: Paciente) : void \\
    + atenderProximo() : Paciente
};

\draw[composicao] (fila.west) -- (paciente.east);

\end{tikzpicture}
\end{center}

\subsection*{Requisitos}
\begin{itemize}
    \item Utilize a classe \texttt{LinkedList} devido à sua eficiência em inserções e remoções no início/meio.
    \item Crie um método \texttt{adicionarPaciente}. Se for uma emergência (prioridade alta), insira o paciente na \textbf{primeira posição} da lista (utilize os métodos específicos do \texttt{LinkedList} ou insira no índice \texttt{0}). Caso contrário, insira no final.
    \item O método \texttt{atenderProximo()} deve remover e retornar o primeiro elemento da lista.
\end{itemize}

% ============================================================
\section*{4) Trabalhando com Set (HashSet e TreeSet)}
Implemente um sistema de registro de CPFs únicos para um evento.

\subsection*{Requisitos}
\begin{itemize}
    \item Primeiramente, utilize um \texttt{HashSet<String>} para tentar inserir diversos CPFs, incluindo alguns duplicados.
    \item Imprima a coleção e comprove que os duplicados foram ignorados e que a ordem de impressão pode ser aleatória.
    \item Em seguida, transfira todos os elementos do \texttt{HashSet} para um \texttt{TreeSet<String>}.
    \item Imprima o \texttt{TreeSet} e verifique se os CPFs foram automaticamente ordenados em ordem alfabética/numérica natural.
\end{itemize}

% ============================================================
\section*{5) Map: Dicionário Chave-Valor (HashMap)}
Desenvolva um sistema acadêmico orientado a objetos que mapeia números de matrícula para os registros completos dos alunos utilizando a estrutura \texttt{HashMap}.

\begin{center}
\begin{tikzpicture}[node distance=2cm and 1cm, , scale=0.8, transform shape]

\node[uml] (aluno) {
    \textbf{Aluno}
    \nodepart{second}
    - nome : String \\
    - curso : String
    \nodepart{third}
    + Aluno(nome: String, curso: String) \\
    + getNome() : String \\
    + getCurso() : String \\
    + toString() : String
};

\node[uml, right=of aluno] (sistema) {
    \textbf{SistemaAcademico}
    \nodepart{second}
    - alunos : Map<Integer, Aluno>
    \nodepart{third}
    + matricularAluno(matricula: int, a: Aluno) : void \\
    + buscarAluno(matricula: int) : Aluno \\
    + listarTodosAlunos() : void
};

\draw[composicao] (sistema.west) -- (aluno.east);

\end{tikzpicture}
\end{center}

\subsection*{Requisitos}
\begin{itemize}
    \item Crie a classe \texttt{Aluno} com os atributos \texttt{nome} e \texttt{curso}, encapsulando-os corretamente e sobrescrevendo o método \texttt{toString()}.
    \item Na classe \texttt{SistemaAcademico}, instancie o atributo \texttt{alunos} utilizando a interface \texttt{Map} e a implementação \texttt{HashMap<Integer, Aluno>}. A chave será o número de matrícula (\texttt{Integer}).
    \item O método \texttt{matricularAluno} deve adicionar o par chave-valor usando \texttt{put()}. Se uma matrícula já existente for passada, permita a sobrescrita mas informe no console que o aluno foi atualizado.
    \item O método \texttt{buscarAluno} deve recuperar e retornar o aluno através do método \texttt{get(Object key)}, ou exibir uma mensagem se a matrícula não existir.
    \item O método \texttt{listarTodosAlunos} deve percorrer a coleção imprimindo as matrículas e os respectivos dados dos alunos. (Dica: utilize \texttt{keySet()} ou \texttt{entrySet()}).
    \item Na classe \texttt{Main}, instancie o sistema, matricule pelo menos 3 alunos, busque por uma matrícula específica e, em seguida, liste todos.
\end{itemize}

% ============================================================
\section*{6) Ordenação Customizada com Comparable}
Modele uma classe \texttt{Produto} e utilize a classe utilitária \texttt{Collections} para ordená-los.

\begin{center}
\begin{tikzpicture}[node distance=2cm and 1cm]

\node[uml, rectangle split parts=2, fill=univSoftBlue!40] (comparable) {
    \textbf{\guillemotleft interface\guillemotright} \\
    \textbf{Comparable<Produto>}
    \nodepart{second}
    + compareTo(outro: Produto) : int
};

\node[uml, below=1.5cm of comparable] (produto) {
    \textbf{Produto}
    \nodepart{second}
    - nome : String \\
    - preco : double
    \nodepart{third}
    + getNome() : String \\
    + getPreco() : double \\
    + toString() : String
};

\draw[implementacao] (produto.north) -- (comparable.south);

\end{tikzpicture}
\end{center}

\subsection*{Requisitos}
\begin{itemize}
    \item A classe \texttt{Produto} deve implementar a interface nativa do Java \texttt{Comparable<Produto>}.
    \item Sobrescreva o método \texttt{compareTo} para definir a ordem natural dos produtos baseada no \textbf{preço} (do menor para o maior).
    \item Crie um \texttt{ArrayList} de produtos e adicione vários objetos de forma desordenada.
    \item Utilize \texttt{Collections.sort(listaDeProdutos)} para ordenar a lista.
    \item Imprima a lista antes e depois da ordenação para validar a implementação. Dica: Sobrescreva o método \texttt{toString()} na classe \texttt{Produto} para facilitar a visualização.
\end{itemize}

\end{document}
